\chapter*{全书结构}

本册分为五个部分，围绕一条主线展开：一个很小的日志统计任务，怎样在现代 CPU、操作系统和网络边界下逐步成长为可测、可并行、可恢复、可分布式的计算系统。第一册已经讲过源码、进程、虚拟地址、汇编、ABI、工具链和基础测量，本册只在建立新成本模型时短暂回扣这些内容。

第一部分是硬件执行模型。它回答一个朴素问题：为什么同样一段循环，在不同输入、布局和线程放置下会表现完全不同。取指前端、I-cache、ITLB、核心流水线、乱序执行、分支预测、寄存器重命名、执行端口、load/store queue、cache、TLB、page cache、NUMA、互连和缓存一致性，共同把源码、机器码、地址序列和线程位置连接到硬件等待。

第二部分是单机高性能计算。它不把优化写成技巧表，而是把性能结论固定为可证伪的论证：输入是什么，reference 是什么，字节数和算术强度是什么，工作集跨过哪些缓存和页边界，编译器生成了什么形状，计数器或替代证据支持哪种瓶颈。数据布局、局部性、预取、SIMD、指令级并行和典型内核，都从 Compute Systems Engine 的热路径中生长出来。

第三部分是多线程、同步与内存模型。它的第一原则是正确性先于性能。线程、调度、亲和性、锁、条件变量、信号量、原子操作、内存序、false sharing、无锁队列和内存回收，都围绕共享状态的三个问题展开：谁能读写，怎样等待，何时可见。若不变量、等待谓词、发布关系、线性化点和对象生命周期没有说清楚，所谓高性能只是偶然没有失败。

第四部分是并行算法与运行时。归约、扫描、分区、任务图、线程池、工作窃取、异步 I/O 和背压，是把多核硬件变成可用计算能力的桥梁。这里的重点不是多写几个 worker，而是让算法依赖、输出位置、任务状态、取消、关闭、buffer 生命周期、completion 和业务提交都能被审查。

第五部分是分布式计算。它把单机原则扩展到网络和集群：消息身份、序列化、deadline、重试、attempt、epoch、分片、复制、共识、shuffle、checkpoint、manifest 和观测性。分布式不是“多几台机器”，而是把通信、迟到响应和部分失败纳入计算模型。单机中的提交事实，在这里会长成跨节点的一致性边界。

\begin{keyidea}
第二册的主线是一条从硬件到集群的因果链：硬件决定局部成本，共享内存决定协作边界，运行时决定任务怎样推进，I/O 决定结果何时完成，分布式系统决定失败后谁是权威事实。
\end{keyidea}

\section*{原理卷与代码卷}

本册正文是原理卷。它使用伪代码、状态表、机制图、实验矩阵和关键 C++ 片段来解释问题，不把完整工程实现塞入每章。完整代码单独进入配套代码实践卷。代码卷负责实现 Compute Systems Engine：输入生成、顺序 reference、parser、HotBatch、benchmark harness、thread-local histogram、stable filter、partition manifest、线程池、有界队列、AtomicProtocolNote、SPSC/MPMC 候选、任务图、可靠写、MessageBus、Checkpoint、RunReport 和故障注入。

这种分离让原理卷可以保持连续叙述，也让代码卷可以按工程标准增长。原理卷每章仍然必须说明本章机制怎样落到 Compute Systems Engine；代码卷每章则给出完整 C++ 实现、测试、测量命令和报告样例。两卷互相引用，但不互相替代。

\section*{从第一册到第二册}

第一册已经讲清楚的内容，本册只作为前提使用：源码到目标文件和 ELF 的路径，进程和虚拟地址空间，x86-64 汇编阅读，System V AMD64 ABI，基础 Linux 命令，基础 benchmark 设计，基础 \cmd{perf stat} 和 Linux 源码入口。第二册遇到这些内容时，重点不是重新定义，而是解释它们在可扩展计算中的新作用。

例如，第一册的虚拟地址空间重点是映射和缺页；第二册的 TLB 和页级性能进一步解释页大小、page walk、NUMA 放置和首次触页怎样改变吞吐。第一册的汇编控制流重点是读懂分支和跳转；第二册的分支预测进一步解释输入分布、预测器状态和流水线回滚怎样限制单核吞吐。第一册的基础 benchmark 重点是避免明显错误；第二册的测量进一步用扩展曲线、PMU、trace、等待时间和故障注入支持工程判断。

\section*{章节质量标准}

每章围绕六件事展开。第一，从一个具体问题或反例开始，避免概念空降。第二，给出朴素方案，明确最简单的正确程序是什么。第三，说明朴素方案在哪种输入、规模、调度或失败下暴露问题。第四，引入机制并解释底层原理。第五，用伪代码、关键片段、实验矩阵或源码入口形成证据。第六，说明边界：什么时候不该使用这个机制，或者它把复杂性转移到了哪里。

本册不接受重复定义、口号化优化、只列 API、只列系统名或模板化实验报告。若一个段落不能帮助现代计算系统的机制、判断或实践变得更清楚，就应删除或并入真正的推导。反过来，若一个难点需要更多篇幅才能讲清，应补充具体输入、边界条件、失败路径、状态表、实验设计和工业案例，而不是用抽象形容词堆深度。

\section*{工业案例的用法}

本册会引用 Linux、glibc、LLVM、oneTBB、OpenMP runtime、DuckDB、ClickHouse、Redis、Kafka、Spark、Flink、RocksDB、Raft 等成熟系统，但工业案例不能替代原理推导。每个案例至少回答五个问题：它面对的真实约束是什么；核心设计原理是什么；它牺牲了什么换取什么；它适合哪些输入、硬件和故障模型；它怎样在 Compute Systems Engine 中复现一个缩小版。

读工业系统时，不要写“某系统使用某机制，所以很好”。要写清机制购买了什么能力。例如 selection vector 购买的是少搬冷字段和延迟材料化的能力，但可能引入间接访问；work stealing 购买的是负载均衡，但可能伤害局部性；manifest 购买的是可恢复提交事实，但会引入提交协议、清理策略和恢复扫描成本。

\section*{全书出口}

原理卷的出口是一套可迁移的分析方法：先定义结果和 reference，再画数据路径、地址路径和状态路径；然后判断瓶颈可能在前端、后端、内存、同步、调度、I/O、网络还是提交协议；最后设计实验、记录证据、承认边界并选择工程方案。配套代码卷把这套方法落实到一个可运行的 Compute Systems Engine 中。
